\section{Conclusion}
\label{sec:conclusion}

We presented SmartPlate, an end-to-end pipeline that converts a single RGB photograph of a plate into a pixel-level food segmentation, a 7-group nutritional breakdown, and a $0$--$100$ healthy-plate score grounded in the Harvard Healthy Eating Plate rubric.
The system fine-tunes a compact SegFormer-B0 backbone (3.74 M parameters) on FoodSeg103 to achieve a $7$-group mIoU of $51.3\%$, then passes the pixel-area proportions through an interpretable, closed-form scoring formula.
A ResNet-50 multi-label baseline confirms that image-level classification can detect food-group presence (mean F1 $= 0.784$) but is structurally incapable of producing the proportional estimates required for dietary scoring, motivating the segmentation approach.

\paragraph{Limitations.}
Several assumptions constrain the current system.
First, the $103$-class mIoU of $14.7\%$ reflects the extreme long tail of FoodSeg103: rare classes with few training images contribute near-zero IoU, and the model generalizes poorly to foods outside the benchmark distribution.
Second, macro estimation assumes a fixed plate mass of $400\,$g; without depth information, per-item weight cannot be inferred from a single image, and the reported macro values are best treated as order-of-magnitude estimates.
Third, the class-to-group mapping was constructed manually and reflects Western nutritional conventions; dishes from other culinary traditions (e.g.\ fermented foods, regional preparations) may be assigned to suboptimal groups.
Fourth, the scoring formula is calibrated to the Harvard rubric, which is itself a simplification of nutritional science; it does not account for individual dietary needs, cooking method, or micronutrient content.

\paragraph{Future work.}
The pipeline's modular design makes a number of extensions tractable.
Replacing SegFormer-B0 with a larger backbone (B2 or B5) or using a promptable segmentation model such as FoodSAM~\cite{lan2023foodsam} could substantially improve fine-grained class accuracy.
Incorporating depth or multi-view imagery (as in Nutrition5k~\cite{thames2021nutrition5k}) would enable plate-mass estimation and more accurate macro prediction.
On the scoring side, the fixed rubric could be replaced by a personalized model that adapts targets to a user's dietary profile, allergies, or clinical goals.
Finally, extending the pipeline to video frames from a meal in progress would enable real-time dietary monitoring beyond the single-image assumption of the current system.
